% Section 09: methodology and scientific status.
% Built only from claims in research/ledgers/09-methodology.md (C-09-xxx).
% Cross-references to other sections use ASSUMED labels, which the integrator
% must reconcile with the labels those sections actually define:
%   sec:argument (01, Bostrom and critics), sec:lattice (03), sec:cosmicrays (04),
%   sec:holographic (05), sec:complexity (06), sec:render (07), sec:infophysics (08).
% Requires amssymb (\gtrsim).

\section{The scientific status of the simulation hypothesis}
\label{sec:methodology}

This section sets out the methodological framework for the rest of the review.
We ask what it would mean to test the claim that the observable universe is a
computer simulation, which versions of the claim can be tested, what a
positive or null result would show, and how published results have been
described to the public. Our conclusion is that only restricted
sub-hypotheses about how a simulator is built can be tested. For every
sub-hypothesis we examine, the proposed signature is also predicted by
non-simulation physics.

\subsection{Standards of testability}
\label{sec:methodology:standards}

\paragraph{Falsifiability.}
On Popper's demarcation criterion, a theory is unscientific if it is
compatible with all possible observations. This holds whether the theory is
compatible by construction or because it has been modified to accommodate them
\cite{Thornton2026karl}. Corroboration counts only when it comes from a
``risky'' prediction that might have failed \cite{Thornton2026karl}. Popper
also held that a theory rescued from failed predictions by ad hoc hypotheses
degenerates into dogma \cite{Thornton2026karl}. Lakatos's standard objection is
that high-level theories are always protected by a ``protective belt'' of
auxiliary hypotheses. On his account they are abandoned when a research
programme stalls, not after a single critical test \cite{Thornton2026karl}.

\paragraph{Underdetermination.}
Stanford distinguishes two forms of underdetermination
\cite{Stanford2023underdetermination}. In the holist form, a failed prediction
can be blamed on an auxiliary assumption rather than on the hypothesis under
test. In the contrastive form, rival theories are equally well supported by
the same evidence. In the limiting case of empirically equivalent theories, no
possible evidence can favour one over the other
\cite{Stanford2023underdetermination}. Descartes' evil demon is the classic
sceptical instance: all sensory experience would be the same whether it came
from the demon or from an external world \cite{Stanford2023underdetermination}.

\paragraph{Bayesian confirmation.}
Evidence $e$ confirms a hypothesis $h$ relative to background $k$ if and only if
$P(h\mid e\wedge k)>P(h\mid k)$. It is neutral if the two are equal
\cite{Crupi2025confirmation}. One standard measure of how far $e$ favours $h$
over its negation is the likelihood ratio
$P(e\mid h\wedge k)/P(e\mid \neg h\wedge k)$ \cite{Crupi2025confirmation}.
We use this measure below. A signature discriminates between a simulation
sub-hypothesis and a non-simulation rival only if their likelihoods for it
differ.

\paragraph{Graded falsifiability and non-empirical assessment.}
Carroll argues that falsifiability tracks two virtues, definiteness and
empiricism. On definiteness, a theory should not be compatible with ``any
conceivable set of facts'' \cite{Carroll2019beyond}. He ranks theories from
category~1, where no conceivable test could return an incompatible result, to
category~5, where there are doable and definitive tests. Category~4 covers
theories that can be tested only over part of their parameter space
\cite{Carroll2019beyond}. On his view the important divide lies between
category~1 and the rest. Category-1 theories ``don't really even have a chance
of being true''. Ideas in category~4 are tested where possible, and interest in
them fades as they are squeezed into smaller regions of parameter space
\cite{Carroll2019beyond}.

A separate debate concerns whether theories may be assessed without empirical
data. Dawid argues that non-empirical assessment plays an important part in
high-energy physics \cite{Dawid2013string}. He identifies three arguments: the
no-alternatives argument, the meta-inductive argument from past success, and
the argument from unexpected explanatory interconnections. He casts all three
as Bayesian confirmation \cite{Dawid2019significance}. Ellis and Silk argued
that exempting speculative theories from experimental verification undermines
science \cite{Ellis2014scientific}. Rovelli argues that non-empirical evidence
cannot move a theory from ``maybe'' to ``reliable''. He adds that
no-alternatives arguments hold only under assumptions that may be false
\cite{Rovelli2019dangers}.

The simulation-hypothesis literature is largely non-empirical in Dawid's
sense. The arguments for it are probabilistic
(Sec.~\ref{sec:argument}), and several arguments against it are a priori or
concern in-principle cost (Sec.~\ref{sec:complexity}). The no-alternatives
argument is plainly unavailable to proponents, because each proposed signature
has a non-simulation explanation (Sec.~\ref{sec:methodology:subhyp}).

\subsection{The unrestricted hypothesis}
\label{sec:methodology:unrestricted}

Let $H_{\rm SIM}$ denote the unrestricted claim that the observable universe
is the output of a computation run by a system outside it. $H_{\rm SIM}$
places no constraint on the simulator's hardware, algorithm, resources or
intentions. Chalmers argues that nothing could prove we are not in a
simulation, ``because any such evidence could be simulated''
\cite{Chalmers2022can}. This passage is a publisher-authorized excerpt of
Ref.~\cite{Chalmers2022reality}. Wolpert gives a precise example of implementation-level
equivalence. Beings simulated through fully homomorphic encryption could not
tell their situation apart from a direct simulation of their laws of physics
\cite{Wolpert2025what}. A perfect simulation of a physical theory $T$ and $T$
itself are therefore empirically equivalent in the sense of
Ref.~\cite{Stanford2023underdetermination}. Their likelihood ratio is unity for
every possible observation. $H_{\rm SIM}$ is thus an underdetermination of the
sceptical, Cartesian kind rather than an ordinary scientific one. In Carroll's
scheme it belongs in category~1. This placement is our application; Carroll
discusses the multiverse, not simulation.

Two qualifications matter. First, the hypothesis is not symmetric under
evidence. Chalmers holds that we could obtain ``very strong evidence'' that we
are simulated, for instance if the simulators showed us the source code, even
though this would fall short of proof \cite{Chalmers2022can}. Hsu and Zee
discuss a physical channel for such a disclosure by a creator: a message
encoded in the cosmic microwave background through the choice of fundamental
Lagrangian \cite{Hsu2006message}. $H_{\rm SIM}$ can therefore in principle be confirmed
by a cooperative simulator, but it cannot be refuted. Second, Bostrom himself
does not claim empirical content for the unrestricted hypothesis. He writes
that the best guide to how simulators set up our world is ``the standard
empirical study of the universe we see''. He also writes that the chief
empirical importance of the simulation disjunct lies in its role in the
trilemma \cite{Bostrom2003are}. The probability assigned to $H_{\rm SIM}$
therefore comes from the prior, supplied by the simulation argument
(Sec.~\ref{sec:argument}), and not from data.

\subsection{Restricted sub-hypotheses and their degenerate rivals}
\label{sec:methodology:subhyp}

A test becomes possible only when $H_{\rm SIM}$ is conjoined with auxiliary
assumptions $A$ about the simulator. Beane, Davoudi and Savage make these
assumptions explicit. They assume a classical computer, a hyper-cubic lattice
and an unimproved Wilson action. From the cosmic-ray spectrum they derive
$b^{-1}\gtrsim 10^{11}$~GeV, and they propose a search for lattice-aligned
anisotropy in the highest-energy cosmic rays \cite{Beane2014constraints}.
Table~\ref{tab:subhyp} classifies the restricted hypotheses discussed in this
review. For each class it lists what would count as evidence and which
non-simulation physics predicts the same signature.

\begin{table}[t]
\centering
\footnotesize
\caption{Restricted simulation sub-hypotheses, their proposed signatures, and
non-simulation physics that predicts the same signature. The status column is
a summary only; see the sections cited for details.}
\label{tab:subhyp}
\begin{tabular}{p{0.14\textwidth}p{0.23\textwidth}p{0.27\textwidth}p{0.20\textwidth}}
\hline
Sub-hypothesis & Signature that would count as evidence & Degenerate non-simulation explanation & Status (summary) \\
\hline
H-grid: fixed lattice & Lattice-specific dispersion and energy cut-off; cubic-symmetric anisotropy of ultra-high-energy cosmic rays \cite{Beane2014constraints} & Lorentz violation in quantum-gravity models \cite{Mattingly2005modern}; minimal-length scenarios \cite{Hossenfelder2013minimal}; Lorentz-invariant causal-set discreteness, which also removes the anisotropy \cite{Dowker2004quantum} & Bounds on specific designs only; see Secs.~\ref{sec:lattice}, \ref{sec:cosmicrays} \\
H-noise: finite resolution & Correlated Planckian position noise in co-located interferometers & Planckian quantum geometry with holographic degrees-of-freedom counting \cite{Hogan2012interferometers} & One model excluded at $4.6\sigma$ incl.\ calibration uncertainty \cite{Chou2016first}; further shear-noise limits \cite{Chou2017interferometric}; see Sec.~\ref{sec:holographic} \\
H-lazy: render on demand & Interference outcomes that depend on the availability of which-path information to an observer \cite{Campbell2017on} & Collapse models (GRW/CSL) \cite{Bassi2013models}; von~Neumann--Wigner observer postulate, as noted by the proposers \cite{Campbell2017on} & Proposals; see Sec.~\ref{sec:render} \\
H-budget: bounded compute & Physics that a bounded simulator cannot compute exactly is absent or approximated & Intractability as a physical principle \cite{Aaronson2005np}; physical limits on computation \cite{Lloyd2002computational} & In-principle arguments, no empirical test; see Sec.~\ref{sec:complexity} \\
H-info: information-physical & New information-thermodynamic laws or effects & Information physics within standard physics \cite{Lloyd2002computational} & Contested; see Sec.~\ref{sec:infophysics} \\
H-patch: adaptive simulator & Glitches; slow drift of constants \cite{Barrow2007living} & Varying-constant theories with a light field coupled to matter \cite{Uzan2011varying} & Constancy tests exist but are not simulation tests \cite{Uzan2011varying} \\
\hline
\end{tabular}
\end{table}

Two general points follow from the table. The first is degeneracy. Suppose a
sub-hypothesis $H_{\rm SIM}\wedge A$ and a physical model $P$ both predict a
signature $s$. Then the likelihood ratio between them for $s$ is close to
unity. Observing $s$ would confirm their disjunction relative to standard
physics, but it would not favour simulation over $P$
\cite{Crupi2025confirmation}. A detected lattice anisotropy would show that
spacetime is discrete in a particular way. A detected deviation from
linear quantum dynamics would support a collapse-type modification. In
neither case would the result show that the discreteness or the collapse was
computed by someone. For a simulation-specific reading, a signature would need
features that the rival models do not motivate. An example is a
discreteness scale far below the Planck scale, which Beane et al.\ note is
possible in the simulation scenario \cite{Beane2014constraints}. Even then,
contrastive underdetermination allows a non-simulation model to be built that
reproduces the feature \cite{Stanford2023underdetermination}.

The second point concerns parameter space. The sub-hypotheses sit in Carroll's
category~4. Beane et al.\ note that lattice improvement ``masks much of our
ability to probe'' the scenario. Apart from dispersion-relation effects,
$\mathcal{O}(b^2)$ operators easily avoid obvious probes
\cite{Beane2014constraints}. A null result therefore excludes one simulator
design at one resolution. A finer or improved lattice always survives. The
same holds for naive preferred-frame discreteness whether or not it is
simulated. Collins et al.\ show that it generically produces percent-level
low-energy Lorentz violation unless it is strongly fine-tuned
\cite{Collins2004lorentz}. Constraints of this kind bear on the physics of
discreteness. They bear on simulation only through the assumed design.

\subsection{Adaptive simulators and the meaning of null results}
\label{sec:methodology:adaptive}

The strongest obstacle to testing is a simulator that responds to being
probed. Bostrom describes this case explicitly. A simulator tracking
observers' belief states could fill in microscopic detail ``on an as-needed
basis''. If an error occurred, it could ``edit the states of any brains that
have become aware of an anomaly'' or ``skip back a few seconds and rerun the
simulation'' \cite{Bostrom2003are}. Barrow argues that simulators with
incomplete knowledge of physics would have to ``patch up the problems one by
one''. Inhabitants would then see occasional glitches and slow drifts in the
constants \cite{Barrow2007living}. Campbell et al.\ model a rendering engine
that must preserve consistency while avoiding detection. They propose
experiments intended to exploit the tension between these two requirements
\cite{Campbell2017on}. Hossenfelder questions whether adaptation is coherent.
She asks how a programmer could notice that a simulated mind is about to
detect a contradiction, and repair it in time, without predicting in advance
what the mind will investigate \cite{Hossenfelder2017no}.

Methodologically, adaptivity acts as an auxiliary hypothesis of the holist
kind \cite{Stanford2023underdetermination}. Suppose the simulator suppresses
detectable anomalies. Then $P(\text{null}\mid H_{\rm SIM}\wedge\text{adaptive})
\approx P(\text{null}\mid\neg H_{\rm SIM})$, and a null result is
approximately relevance-neutral \cite{Crupi2025confirmation}. Invoking adaptivity after a
null result has the structure that Popper classed as an ad hoc rescue
\cite{Thornton2026karl}. Rovelli's point applies here too. A flexible theory
loses Bayesian credibility with each failed prediction only when a success
would have raised it \cite{Rovelli2019dangers}. For a fully adaptive simulator
no outcome was expected, so null results cost it nothing. Barrow's patching
scenario does yield a signature, drifting constants. That signature is itself
degenerate with varying-constant physics, in which a varying constant reflects
a light field coupled to matter \cite{Uzan2011varying}.

We therefore state every null result in this review as conditional. It
constrains simulators that are non-adaptive (they do not intervene in response
to probes) and that implement the assumed architecture. The finite-resource
premise does real work here. Beane et al.\ argue that finite simulator
resources imply a non-zero lattice spacing, so that ``in principle there
always remains the possibility for the simulated to discover the simulators''
\cite{Beane2014constraints}. That conclusion holds only under the finite-resource
premise.

\subsection{Published assessments of scientific status}
\label{sec:methodology:assessments}

Published assessments range from outright rejection of the hypothesis's
scientific standing to the view that it cannot be ruled out.
Hossenfelder calls the hypothesis unscientific, ``not a serious scientific
argument'', while stating that this ``doesn't mean it's wrong''
\cite{Hossenfelder2021simulation}. Her central objection is that it assumes,
without explanation, that all observations can be reproduced by an algorithm
other than the confirmed laws. She adds that algorithmic reproductions of
natural laws are usually incompatible with relativistic symmetries
\cite{Hossenfelder2021simulation}. Aaronson argues that the Ringel--Kovrizhin result
(Sec.~\ref{sec:methodology:reception}) could not rule out the hypothesis. Among
his reasons are that the simulator could be a quantum computer, or a classical
one given exponentially more time \cite{Aaronson2017because}. Vazza derives energy requirements that he
finds incompatible with physics for any simulator in a universe like ours
\cite{Vazza2025astrophysical}. He concludes that the hypothesis can be
``reasonably well tested only with respect to universes which are at least
playing according to the Physics play book''. Anything else, he writes, lies
``beyond the bounds of falsifiability'' \cite{Vazza2025astrophysical}. Wolpert
argues that experimental approaches can give only partial answers, because
they rest on assumptions such as a digital simulating computer
\cite{Wolpert2025what}. Chalmers argues that we should take the hypothesis
seriously and cannot rule it out \cite{Chalmers2024taking}. His epistemic
asymmetry is discussed above (see also Sec.~\ref{sec:argument}). The
hypothesis was the subject of the 2016 Isaac Asimov Memorial Debate at the
American Museum of Natural History, with physicists and a philosopher on the
panel \cite{AMNH2016is}. We cite the debate as an event only.

Among peer-reviewed work, the explicit discussions of testability that we
found are physics or computer-science papers, not philosophy of science
\cite{Beane2014constraints,Campbell2017on,Vazza2025astrophysical,Wolpert2025what}.
We are not aware of a peer-reviewed philosophy-of-science article devoted
specifically to the falsifiability of the simulation hypothesis. Our search is
described in the accompanying research notes.

\subsection{Reception: what was reported versus what was claimed}
\label{sec:methodology:reception}

Table~\ref{tab:reception} pairs four results with how they were publicly
reported. The news items and press releases cited here are evidence of
reporting only, not of any scientific claim.

\begin{table}[t]
\centering
\footnotesize
\caption{Claims made in four papers compared with how the results were
reported. News and press sources document reporting only.}
\label{tab:reception}
\begin{tabular}{p{0.16\textwidth}p{0.38\textwidth}p{0.34\textwidth}}
\hline
Paper & What the paper claims & How it was reported \\
\hline
Beane et al.\ \cite{Beane2014constraints} & Bounds and a proposed signature for one assumed design (unimproved Wilson action, cubic lattice); improvement ``masks'' the signal & University release: ``Researchers say idea can be tested''; ``the first testable signature of such an idea'' \cite{UW2012do} \\
Ringel \& Kovrizhin \cite{Ringel2017quantized} & Quantized gravitational responses obstruct local sign-free quantum Monte Carlo; no mention of the universe or the simulation hypothesis & ``Physicists Confirm That We're Not Living In a Computer Simulation'' \cite{Eck2017physicists} \\
Vopson \cite{Vopson2023second} & Abstract: ``scientific evidence that appears to underpin the simulated universe hypothesis'' & ``\ldots Could Mean We Really Live in a Simulation''; author quoted: the law ``is valid regardless of whether the universe is a simulation or not'' \cite{Ferreira2023mind} \\
Faizal et al.\ \cite{Faizal2025consequences} & Incompleteness theorems applied to quantum gravity modelled as an axiomatic, algorithmic structure: ``the universe cannot be a simulation'' & ``mathematically proven'' and ``a definitive answer'' \cite{UBCO2025ubco}; ``Physicists prove the Universe isn't a simulation after all'' \cite{ScienceDaily2025physicists} \\
\hline
\end{tabular}
\end{table}

The four cases differ in kind. In the Beane et al.\ case the release is
broadly faithful. It nevertheless elides the dependence on one assumed design,
which the paper itself stresses \cite{Beane2014constraints,UW2012do}. In the
Ringel--Kovrizhin case the paper concerns a class of classical algorithms and
never mentions the simulation hypothesis \cite{Ringel2017quantized}. The
simulation conclusion was added in reporting \cite{Eck2017physicists}.
Aaronson responded publicly that the hypothesis ``has not been falsified;
remains unfalsifiable'' \cite{Aaronson2017because}. One outlet reported that
the authors were surprised by the headlines \cite{Galeon2017fact}. In the
Vopson case the paper's abstract is itself hedged (``appears to underpin'')
\cite{Vopson2023second}. The author's statement, quoted in the coverage, that
the proposed law holds whether or not the universe is simulated amounts to a
concession of degeneracy \cite{Ferreira2023mind}. In the Faizal et al.\ case
the strong conclusion appears in the paper itself: the simulation hypothesis
is ``logically impossible rather than merely implausible''
\cite{Faizal2025consequences}. The press release then presented it as a
definitive mathematical proof \cite{UBCO2025ubco}. The argument's validity is
disputed. A comment posted to arXiv argues that undecidability limits
provability, not execution \cite{Redden2025provability}. We assess this in
Sec.~\ref{sec:complexity}. In every case the headline tied the result to
the unrestricted hypothesis (``the idea'', ``we'', ``the Universe''). By the
analysis of Sec.~\ref{sec:methodology:unrestricted}, no single result can
settle that hypothesis.

\subsection{Working thesis of this review}
\label{sec:methodology:thesis}

On this analysis we adopt the following position. The unrestricted
simulation hypothesis is not an empirical hypothesis. No observation can refute
it, and strong evidence for it would most plausibly require disclosure by the
simulator itself \cite{Chalmers2022can}.
What can be tested are conjunctions of the hypothesis with explicit assumptions
about simulator architecture, resources and non-adaptivity. For each such
sub-hypothesis in Table~\ref{tab:subhyp}, the signature is also predicted by
non-simulation physics. A detection would therefore establish new physics,
such as discreteness, holographic noise or collapse. It would not establish
that the universe is simulated. A null result bounds the space of
non-adaptive, resource-limited simulator designs that share the assumed
architecture. Accordingly, in what follows we speak of constraining simulator
designs, and we avoid ``proving'' or ``disproving'' the simulation hypothesis.
